\section{Introduction}

\subsection{ETA as an Operational Context}
Estimated time of arrival (ETA) is a central prediction task in transportation systems, informing navigation, dispatch, logistics, and traveler information services. In such settings, ETA is not merely a descriptive quantity: it shapes route choice, scheduling, service reliability, and user trust. For operational platforms, inaccurate or unstable ETA estimates propagate into missed coordination, inefficient routing, and degraded service quality~\cite{modica2024eta}.

\subsection{Route-Aware ETA Data and Representation Gaps}
The central difficulty addressed in this work is not the absence of traffic data in general, but the lack of open, reusable, route-aware data representations tailored to ETA-oriented tasks. Much of the traffic prediction literature relies on sensor-centric or aggregate spatio-temporal datasets, while trajectory-based data often captures movement without preserving the route-aware structure characteristic of navigation-style ETA settings~\cite{yu2017trafficpred}. This problem is compounded by the scarcity of openly reusable ETA datasets: many operationally valuable sources remain proprietary, making reproducible experimentation difficult.

The gap is therefore twofold. First, ETA-oriented learning requires data in which road topology, traffic evolution, and route context are aligned in one coherent form. Second, that representation should support controlled downstream validation rather than ending at offline dataset export. These requirements become harder still when one seeks to accommodate both simulation-derived traffic data and real trajectory data within a common framework.

\subsection{Two Data Sources, One Shared Representation}
This paper addresses that gap through a common ETA-oriented representation constructed from two fundamentally different data sources. The first path begins with controlled simulation, where the evolving traffic state and planned routes are directly observable. The second begins with real trajectories, where route inference and data repair are required before the data can be expressed in a coherent downstream form.

Although these paths do not produce the same dataset, they produce distinct dataset instances expressed in the same graph-based representation. In this formulation, roads are modeled as edges, while junctions and vehicles are modeled as nodes. The representation is explicitly route-aware: vehicles are tracked along known or inferred planned routes, making the data substantially closer to navigation-system usage than to generic traffic datasets. It is also dynamic in two senses: traffic evolution appears in time-varying node and edge features, and the graph structure itself changes as vehicles move and dynamic traffic relations are created and removed. While designed primarily for vehicle-level ETA-oriented tasks, the same representation can also support other downstream objectives, such as road-level load estimation or area-level traffic analysis.

\subsection{Integrated Artifact: Dataset Tool and Testing Environment}
The concrete contribution of this work is an integrated artifact ecosystem composed of a desktop dataset-construction tool and a web-based testing environment. The desktop tool supports both data-construction paths. For simulation-based datasets, it provides an interactive, user-facing desktop environment for configuring scenario behavior, controlling simulation settings, visually validating key semantic elements such as zones, landmarks, and vehicle classes, and exporting route-aware graph snapshots with adjustable temporal granularity. For real-trajectory datasets, it supports network selection, trajectory repair, route alignment, and conversion into the same ETA-oriented representation.

The web environment complements dataset construction by exercising the produced representation in controlled downstream use. It supports the attachment of a representation-compatible ETA predictor, interactive trip-based testing, registration of prediction outcomes in a database, statistical analysis of trip and model behavior, generation of plots, and export of PDF reports. As an open-source evaluation framework, it can also be reused with other compatible predictors and alternative configurations beyond the specific demonstrator instance described in this paper. Artifact URLs are omitted for double-blind review and can be restored in the camera-ready version.

\subsection{Validation Scope and Model Role}
The predictive model used within the workflow is not the principal contribution of this article. Its purpose is methodological: it serves as a representation-compatible validation component through which the usability of the produced datasets can be assessed. The evidence reported later in the paper should therefore be interpreted as support for representation coherence, trainability, and end-to-end workflow functionality, rather than as a claim of predictive-model novelty or superiority.

This distinction is particularly important because the representation exposes route-aware and graph-structured information not naturally accommodated by many off-the-shelf baselines. The attached predictor is therefore used as a practical downstream instrument for exercising the representation under controlled conditions. Although the primary goal of the web environment is not to deliver a production navigation system, it does provide evidence that the proposed representation is operational enough to support navigation-style ETA testing in an interactive setting.

\subsection{Paper Organization}
The remainder of this paper is organized as follows. The next section positions the work relative to existing ETA datasets, graph-based traffic representations, simulation and trajectory-conversion pipelines, and traffic evaluation environments. The following sections then present the system overview, the dataset-construction paths and shared representation, and the web-based evaluation loop. The paper concludes with an evaluation snapshot, a discussion of limitations and implications, and directions for future work.

